% ============================================================================
% METHODOLOGY - REVISED
% ============================================================================

\section{Methodology}

\subsection{Problem Definition and Core Concepts}

A Constraint Satisfaction Problem is formally defined as $\mathcal{P} = (\mathbf{X}, \mathbf{D}, \mathbf{C})$ where $\mathbf{X} = \{x_1, \ldots, x_n\}$ is a set of variables, $\mathbf{D} = \{D_1, \ldots, D_n\}$ specifies finite domains ($x_i \in D_i$), and $\mathbf{C} = \{c_1, \ldots, c_m\}$ is a set of constraints. A solution is a complete assignment $\theta: \mathbf{X} \to \bigcup_i D_i$ such that $\theta(x_i) \in D_i$ and all constraints are satisfied. We focus on logical grid puzzles (e.g., Zebra puzzles) where variables represent entity-attribute pairs and constraints are expressed in natural language.

\textbf{Solving Trajectory:} For puzzle instance $\mathcal{P}$, a solving trajectory is a temporal sequence of reasoning states and actions:
\[
\mathcal{T} = \langle (S_0, a_1, S_1), (S_1, a_2, S_2), \ldots, (S_{T-1}, a_T, S_T) \rangle
\]
where $S_t = (C_t, \delta_t)$ is the solver state at step $t$, with $C_t \subseteq \mathbf{C}$ being the current set of formalized constraints and $\delta_t: \mathbf{X} \to 2^{\bigcup_i D_i}$ being domain assignments (each variable maps to a set of possible values). Action $a_t$ is an LLM-generated inference: a natural language statement plus its Z3 formalization. A trajectory is \textit{successful} if the final state $S_T$ satisfies all constraints (Z3 returns SAT) and the variable assignment matches ground truth.

\textbf{Solving Paradigm:} We define a solving paradigm as a formally verifiable five-tuple:
\[
P = (\mathcal{Q}, \mathcal{O}, \phi_{\text{pre}}, \phi_{\text{post}}, \mathcal{S})
\]
\begin{itemize}
    \item $\mathcal{Q}$: trigger conditions specifying constraint types and solver state properties that activate the paradigm.
    \item $\mathcal{O}$: inference operation, given as a parameterized template (e.g., $\texttt{position}(X) = k + d$).
    \item $\phi_{\text{pre}}$: Z3-formalized precondition assertion describing properties the constraint set must satisfy before applying the paradigm.
    \item $\phi_{\text{post}}$: Z3-formalized postcondition assertion describing what the paradigm's operation should establish.
    \item $\mathcal{S}$: scope documentation indicating applicable puzzle types and scales where the paradigm has been verified.
\end{itemize}

Critically, $\phi_{\text{pre}}$ and $\phi_{\text{post}}$ are formal logical statements mechanically verifiable by Z3. This distinguishes PRISM from prior memory-based methods that store unverified natural language heuristics.

\subsection{PRISM Architecture Overview}

PRISM operates via two decoupled stages: offline paradigm distillation and online guided reasoning with adaptive repair.

\textbf{Offline Stage:} Given $N$ training puzzles with ground truth solutions, execute the following pipeline (§3.2):
\begin{enumerate}
    \item Trajectory Collection: Run LLM+Z3 system on training puzzles to gather successful solving trajectories.
    \item KDP Identification: Locate Key Decision Points—steps with significant domain reductions or non-trivial inference types.
    \item Cross-Trajectory Clustering: Group similar KDPs using constraint-type features and complete-linkage agglomerative clustering.
    \item Paradigm Abstraction: For each cluster, synthesize a reusable paradigm specification via LLM.
    \item Z3 Triple Verification: Test each paradigm candidate for soundness, effect (non-vacuousness), and trigger precision. Only candidates passing all three checks are admitted.
    \item Library Ingestion: Persist verified paradigms in the Paradigm Library $\mathcal{L}$.
\end{enumerate}
Result: ~35 high-confidence paradigms extracted from ~1500 training trajectories.

\textbf{Online Stage:} For each new puzzle (§3.3--3.4):
\begin{enumerate}
    \item Initial Translation: LLM translates natural language constraints to Z3 formulas.
    \item State-Driven Paradigm Retrieval: Extract constraint-type signature from current solver state; retrieve paradigms via Layer-1 (type matching) and Layer-2 (LLM semantic validation).
    \item Consistency Pre-Check: Verify each candidate paradigm's operation is satisfiable under current constraints before injection.
    \item Hint Injection: Pass verified paradigms as structured prompts to guide LLM inference.
    \item Step Verification: Z3 verifies each LLM-generated inference; returns SAT/UNSAT/ERROR.
    \item Repair Loop: If UNSAT, extract the UNSAT core (minimal contradiction-causing constraint subset), record repair attempt in Repair Trajectory Memory, and detect stagnation/loops.
    \item Strategy Escalation: Trigger four-level escalation if stagnation/loop detected.
    \item Write-Back: Upon success, integrate effective repairs into $\mathcal{L}$.
\end{enumerate}

\subsection{Offline Stage: Paradigm Distillation}

\subsubsection{Trajectory Collection}

Given $N = 600$ training puzzles across difficulty levels (easy/medium/hard), we execute LLM+Z3 solving with generation temperature $T = 0.7$ to induce trajectory diversity. We run $r = 3$ independent generations per puzzle. We retain only successful trajectories where Z3 returns SAT with variable assignments matching ground truth. This yields approximately 1500 successful trajectories.

Each trajectory step $a_t$ carries metadata:
\begin{itemize}
    \item Z3 verification outcome (SAT/UNSAT/ERROR)
    \item Domain change record $\Delta_t$ (variables with reduced domains)
    \item Constraint-type signature of current solver state
    \item Preliminary step-type classification
\end{itemize}

\subsubsection{Key Decision Point Identification}

A step $(S_{t-1}, a_t, S_t)$ in trajectory $\mathcal{T}$ qualifies as a Key Decision Point if it satisfies at least one of:

\textbf{Condition A (Significant Domain Reduction):} $\exists x \in \Delta_t : |\delta_t(x)| \leq 1$. That is, at least one variable's domain is uniquely determined in this step.

\textbf{Condition B (Non-Trivial Inference Type):} The step is classified as CHAIN\_PROPAGATION (transitive reasoning) or CONTRADICTION\_ELIMINATION (explicit elimination of infeasible assignments).

Step types are assigned via lightweight rule-based classification applied to the natural language action text. Across 1500 trajectories, this process extracts approximately 6000 KDPs.

\subsubsection{Cross-Trajectory Clustering and Paradigm Abstraction}

\textbf{Feature Representation:} Each KDP is represented as a fixed-dimensional vector:
\[
\mathbf{f}(\text{kdp}) = [\mathbf{h}_{\text{ct}}, \mathbf{b}_{\text{dom}}, \mathbf{e}_\tau, n_{\text{var}}, n_{\text{con}}]
\]
where:
\begin{itemize}
    \item $\mathbf{h}_{\text{ct}} \in \mathbb{R}^{|\mathcal{T}_c|}$ is a normalized histogram of constraint types in the current solver state. $\mathcal{T}_c$ is a vocabulary of 10 constraint types (direct assignment, adjacency, relative position, exclusion, binding, etc.).
    \item $\mathbf{b}_{\text{dom}} \in \mathbb{R}^4$ is a domain-size distribution histogram (counts variables with domain size 1, 2, 3--5, $\geq 6$).
    \item $\mathbf{e}_\tau$ is a one-hot encoding of the KDP's step type.
    \item $n_{\text{var}}, n_{\text{con}}$ are normalized problem scale metrics (variable and constraint counts, divided by maximum observed values).
\end{itemize}

\textbf{Clustering:} We apply complete-linkage agglomerative clustering using cosine distance. Clusters are formed with a distance threshold $\theta = 0.25$ (determined via grid search on a development set). Clusters with fewer than $m_{\min} = 5$ members are discarded to ensure statistical reliability.

\textbf{Paradigm Abstraction:} For each cluster $C_k$ with at least 5 members, we uniformly sample up to 10 representative KDPs. For each KDP, we format: (1) the solver state (current constraints, domain assignments), (2) the reasoning text from the original trajectory, (3) domain change indicators. We then prompt an LLM:

\begin{quote}
\textit{Given the following similar reasoning steps from successful solving trajectories [formatted examples], extract a general solving paradigm in JSON format with fields: trigger (constraint types and state features), operation (parameterized inference template), precondition\_z3 (Z3-formalized assumptions), postcondition\_z3 (Z3-formalized claims), scope (applicability).}
\end{quote}

The LLM generates a paradigm candidate $\hat{P}_k$, which enters the verification stage. Up to 2 retries per cluster are permitted to account for LLM variability.

\subsubsection{Z3 Triple Verification}

Each candidate paradigm $\hat{P}_k$ undergoes three independent verification checks. \textit{Any failure results in rejection.}

\textbf{(1) Soundness:} The paradigm's operation must be satisfiable under its precondition.
\begin{itemize}
    \item Sample 50 random puzzle solver states satisfying $\hat{P}_k$'s trigger conditions.
    \item For each state $i$, instantiate $\hat{P}_k$'s operation template with the current bindings and test: $\text{Z3-SAT}(C_i \cup \{\text{op}(\hat{P}_k)\})$.
    \item Compute: $\text{soundness}(\hat{P}_k) = \frac{1}{50}\sum_{i=1}^{50} \mathbf{1}[\text{Z3-SAT}(...)]$.
    \item Require: $\text{soundness}(\hat{P}_k) \geq 0.90$.
\end{itemize}

\textbf{(2) Effect:} The paradigm's operation must genuinely constrain the solution space—it cannot be vacuous or already implied by the precondition.
\begin{itemize}
    \item On the 50 states from Soundness check, verify that applying the operation causes net domain reduction: count variables with reduced domains.
    \item Compute: $\text{effect}(\hat{P}_k) = \frac{\text{states with domain reduction}}{50}$.
    \item Require: $\text{effect}(\hat{P}_k) \geq 0.80$.
\end{itemize}

\textbf{(3) Trigger Precision:} The paradigm's trigger conditions must be selective—they should not fire on inappropriate states.
\begin{itemize}
    \item Sample 20 random solver states that do \textit{not} satisfy $\hat{P}_k$'s trigger conditions.
    \item Test: $\text{Z3-UNSAT}(C_i \cup \{\text{op}(\hat{P}_k)\})$ for each state (the operation should be unsatisfiable).
    \item Compute: $\text{precision}(\hat{P}_k) = \frac{1}{20}\sum_{i=1}^{20} \mathbf{1}[\text{Z3-UNSAT}(...)]$.
    \item Require: $\text{precision}(\hat{P}_k) \geq 0.80$.
\end{itemize}

Paradigms passing all three checks are assigned a confidence score and persisted in the SQLite-backed Paradigm Library $\mathcal{L}$. Across 600 training puzzles, the offline stage yields approximately 35 verified paradigms, representing a 40:1 compression ratio.

\subsection{Online Stage: Paradigm-Guided Reasoning and Adaptive Repair}

\subsubsection{Initial Translation and Constraint Formalization}

Given a new puzzle $\mathcal{P}^*$ in natural language, the LLM executes an initial translation to convert constraints into Z3 formulas. The translator outputs a set of Z3 constraints $C^* = \{c_1, \ldots, c_m\}$ and initializes a solver state $S_0 = (C_0, \delta_0)$ where $\delta_0$ assigns each variable its full domain.

\subsubsection{Two-Layer Paradigm Retrieval}

At each reasoning step $t$, we extract the current constraint-type signature $\sigma(S_t) = \text{sorted}(\{type(c) : c \in C_t\})$ and retrieve relevant paradigms using two filtering layers:

\textbf{Layer 1 (Type Matching):} Return all paradigms whose trigger constraint-type requirements are subsumed by the current signature:
\[
\mathcal{L}_1(S_t) = \{P \in \mathcal{L} : \mathcal{Q}_P.\text{required\_types} \subseteq \sigma(S_t) \land \text{conf}(P) \geq \tau_{\min}\}
\]
This layer executes in $O(|\mathcal{L}|)$ time and typically reduces candidates from ~35 to 3--5.

\textbf{Layer 2 (LLM Semantic Judgment):} If $|\mathcal{L}_1(S_t)| \geq 2$, construct a prompt asking the LLM to judge which paradigms have trigger conditions genuinely satisfied by the current state:

\begin{quote}
\textit{Given current constraints and solver state [formatted description], which of these paradigm triggers apply? (y/n for each)}
\end{quote}

This layer is invoked at most once per step and introduces at most one additional LLM call. Return the top $k = 3$ paradigms ordered by confidence, denoted $\mathcal{L}_2(S_t)$.

\subsubsection{Consistency Pre-Check and Hint Injection}

Before injecting paradigms as hints, perform a Z3 consistency check for each candidate $P \in \mathcal{L}_2(S_t)$:
\[
\text{consistent}(P, S_t) := \text{Z3-SAT}(C_t \cup \{\text{op}(P, \cdot)\})
\]

If the paradigm's operation is unsatisfiable under current constraints, remove it from the injection set. This acts as a second safety barrier, ensuring suggested paradigms are logically feasible.

The final injection set $\mathcal{L}_{\text{inject}}(S_t)$ consists of paradigms passing consistency check. Inject these as structured hints in the LLM prompt:

\begin{quote}
\textit{Current constraint state: [constraints]. Suggested solving patterns: (1) [Paradigm 1 trigger, operation, and rationale]. (2) [Paradigm 2...]}
\end{quote}

\subsubsection{Step Verification and UNSAT Core Extraction}

The LLM generates a proposed reasoning step (natural language + Z3 formalization). Z3 verifies the step:
\begin{itemize}
    \item \textbf{SAT:} Accept the step, update solver state $S_{t+1}$, and continue.
    \item \textbf{UNSAT:} Extract the UNSAT core—the minimal subset of constraints $U_t \subseteq C_t$ such that $U_t \cup \{\text{proposed\_operation}\}$ is unsatisfiable. Record the failure for repair analysis.
    \item \textbf{ERROR:} Parse failure message, attempt lower-level correction.
\end{itemize}

\subsubsection{Repair Trajectory Memory and Adaptive Strategy Escalation}

When step verification fails (UNSAT), record the event in Repair Trajectory Memory:
\[
R_t = (\text{error\_type}, U_t, \text{repair\_action}, \text{outcome})
\]

The error type is inferred from UNSAT core analysis (over-constrained translation, missing binding, contradiction in domain assignments, etc.). The memory maintains a circular buffer of recent repair attempts.

\textbf{Stagnation Detection:} Compute Jaccard similarity between consecutive UNSAT cores:
\[
J_t = \frac{|U_t \cap U_{t-1}|}{|U_t \cup U_{t-1}|}
\]

If $J_t \geq \tau_{\text{stag}} = 0.75$ for three consecutive repair rounds, stagnation is detected (the system is repeatedly hitting the same constraints).

\textbf{Loop Detection:} Embed each repair action $r_t$ into a feature vector (repair type, target constraints, magnitude of change). Compute cosine similarity between recent actions:
\[
\cos_t = \frac{\mathbf{r}_t \cdot \mathbf{r}_{t-k}}{|\mathbf{r}_t| |\mathbf{r}_{t-k}|}
\]

If $\cos_t \geq \tau_{\text{loop}} = 0.90$ for repeated repair attempts with different indices, loop is detected (same action repeated).

\textbf{Four-Level Strategy Escalation:}
\begin{enumerate}
    \item \textbf{L1 (Switch Repair Target):} If stagnation detected, switch to repairing a different constraint from the UNSAT core.
    \item \textbf{L2 (Switch Repair Type):} If L1 fails, switch the type of repair (e.g., from retraction to constraint relaxation).
    \item \textbf{L3 (Revert to SAT Checkpoint):} If L2 fails, revert to the most recent solver state that returned SAT and take a different inference path.
    \item \textbf{L4 (Full Retranslation):} If all above fail, discard current constraint formalization and translate the entire puzzle from scratch with modified prompting.
\end{enumerate}

Each level triggers after at most $K = 8$ repair rounds at the previous level.

\subsubsection{Write-Back and Library Augmentation}

Upon successful solving of puzzle $\mathcal{P}^*$, extract effective repair patterns from the trajectory and write them back to $\mathcal{L}$:
\begin{itemize}
    \item For each repair action that led to progress (visible domain reduction or movement toward solution), extract its precondition and postcondition.
    \item If the pattern is novel (not matched by existing paradigm triggers), create a candidate paradigm.
    \item Apply Z3 triple verification (accelerated: lower sampling $n = 20$ to reduce overhead).
    \item If verified, append to $\mathcal{L}$ and update confidence scores.
\end{itemize}

This enables paradigm libraries to grow and adapt over time.
